\section{Robustness, and a mechanism test that could have failed}
\label{app:robustness}

Two questions are asked here, both against the \emph{oracle} Mat\'ern tier rather
than the deployable one, because the interesting failure would be the derived
prior falling behind the best a tuned kernel can do.

\subsection{Does the derived prior stay level as conditions change?}
Three one-dimensional sweeps on the one-wall layout, each varying a single thing
and holding the rest at the main table's values: the observation ratio, the
observation noise, and how wrong the nominal coefficients are.  The last is the
one that matters for the claim, because the whole construction rests on being
allowed to be wrong about the coefficients: the main table uses values wrong by
$50\%$, and the sweep asks how much further that can be pushed before the prior
stops being useful.

\begin{table}[h]
  \centering
  \caption{Sweep over the fraction of entries observed, one-wall layout.  The Mat\'ern column is the oracle tier, tuned against the held-out region.}
  \label{tab:sweep-ratio}
  \small
  \begin{tabular}{lcccc}
    \toprule
    fraction of entries observed & ours & Mat\'ern$^\star$ & neural CP & margin \\
    \midrule
    $0.003$ & 0.5450 & 0.5388 & 0.8723 & $-0.0062$ \\
    $0.005$ & 0.5429 & 0.5390 & 0.8622 & $-0.0040$ \\
    $0.01$ & 0.5457 & 0.5433 & 0.8935 & $-0.0024$ \\
    $0.02$ & 0.5517 & 0.5430 & 0.8883 & $-0.0087$ \\
    $0.05$ & 0.5606 & 0.5526 & 0.8916 & $-0.0080$ \\
    \bottomrule
  \end{tabular}
\end{table}

\begin{table}[h]
  \centering
  \caption{Sweep over the observation noise standard deviation, one-wall layout.  The Mat\'ern column is the oracle tier, tuned against the held-out region.}
  \label{tab:sweep-noise}
  \small
  \begin{tabular}{lcccc}
    \toprule
    observation noise standard deviation & ours & Mat\'ern$^\star$ & neural CP & margin \\
    \midrule
    $0.02$ & 0.5482 & 0.5442 & 0.9013 & $-0.0041$ \\
    $0.05$ & 0.5457 & 0.5433 & 0.8935 & $-0.0024$ \\
    $0.15$ & 0.5491 & 0.5459 & 0.9024 & $-0.0033$ \\
    $0.4$ & 0.5631 & 0.5462 & 0.9065 & $-0.0169$ \\
    \bottomrule
  \end{tabular}
\end{table}

\begin{table}[h]
  \centering
  \caption{Sweep over the nominal coefficients, as a factor of the truth, one-wall layout.  The Mat\'ern column is the oracle tier, tuned against the held-out region.}
  \label{tab:sweep-coefficient}
  \small
  \begin{tabular}{lcccc}
    \toprule
    nominal coefficients, as a factor of the truth & ours & Mat\'ern$^\star$ & neural CP & margin \\
    \midrule
    $0.1$ & 0.6736 & 0.5433 & 0.8935 & $-0.1303$ \\
    $0.3$ & 0.5813 & 0.5433 & 0.8935 & $-0.0380$ \\
    $1$ & 0.5457 & 0.5433 & 0.8935 & $-0.0024$ \\
    $3$ & 0.5360 & 0.5433 & 0.8935 & $+0.0072$ \\
    $10$ & 0.5550 & 0.5433 & 0.8935 & $-0.0117$ \\
    \bottomrule
  \end{tabular}
\end{table}



\paragraph{How to read it.}  We are looking for the derived prior to track the
oracle across the sweep, not to beat it.  A margin near zero throughout is the
success case.  Where the prior loses badly is where the construction should not
be used, and the coefficient sweep is the one that draws that boundary.

\paragraph{What the sweeps say.}  Sample count and noise are both benign: across
a $17\times$ range of observations and a $20\times$ range of noise the derived
prior stays within $3\%$ of the oracle.  Note also what \emph{neither} arm does
-- going from $786$ observations to $13\,107$ leaves the error at one wall
essentially unchanged ($0.545$ to $0.561$).  Under a confined layout the binding
constraint is the geometry, not the sample size, which is why adding sensors to
a wall is not a substitute for a prior.

Coefficient error is the sweep that draws a boundary, and it is strongly
asymmetric.  Overestimating the diffusivities is nearly free -- $3\times$ is
slightly \emph{better} than nominal at $+1.3\%$ and $10\times$ costs $2.2\%$ --
while underestimating is not: $0.3\times$ costs $7\%$ and $0.1\times$ costs
$24\%$.  The physics says why.  Overestimating diffusion makes the prior too
smooth, which merely over-regularises; underestimating it leaves the prior
permitting high-frequency structure the field does not have, and under
extrapolation that structure is invented rather than merely tolerated.  The
usable range on this field is roughly $1\times$ to $10\times$, and the main
table's $1.5\times$ sits inside it.

\subsection{Does the advantage follow the axis, as the mechanism claims?}
The explanation offered for \Cref{tab:layouts} is specific: a per-mode spectrum
states how the field decays along one axis, and a one-wall layout forces
extrapolation along one axis, so the two line up.  A weaker reading -- that
physics is generically useful in hard settings -- predicts the same table.

They differ on \emph{which} wall.  The field is anisotropic, $D_x = 0.02$ against
$D_y = 0.006$, so $x$ is the quickly diffusing smooth axis and $y$ the slow rough
one.  A generic kernel commits to a single length scale and is therefore wrong by
more along $y$.  If the gain comes from knowing per-axis decay, the $y$ wall must
show a larger margin than the $x$ wall; if physics is merely generically helpful,
the two walls should agree.  The refutation condition and an isotropic control --
where $D_x = D_y$ and the two walls \emph{must} agree -- were written into the
script before it was run, and the verdict is computed rather than narrated.

% generated by experiments/make_appendix_tables.py; run pending
\emph{(pending: the run for this table has not finished.)}


\paragraph{The prediction failed, and the control rules out the rescue.}
The $y$ wall does not show a larger margin.  It shows a catastrophic loss: our
arm scores $1.2652$ against $0.8445$ for the oracle Mat\'ern, losing on all five
seeds, while at the $x$ wall the two arms are level.  If anisotropy were the
cause, making the field isotropic should have removed the gap.  It does not --
the isotropic control loses by $36.8\%$ at the $y$ wall and ties at the $x$ wall,
a gap of $-37.2$ points where the mechanism requires approximately zero.  So the
axis-alignment story offered for \Cref{tab:layouts} is not supported, and we
withdraw it as an explanation.  Which wall the sensors sit on matters enormously,
and we do not currently know why.

\paragraph{What the failure did reveal.}
Across every experiment in this paper, exactly one arm ever scores worse than
predicting the mean, and it is always ours: $1.0212$ at the corner patch,
$1.2652$ and $1.2240$ at the two $y$-wall conditions.  No baseline ever exceeds
$0.95$.  This is the more useful finding, and it is not about anisotropy.  A
generic kernel with a long length scale degrades gracefully: when the geometry
defeats it, it shrinks toward the mean and stops there, which is why the oracle's
chosen length scale at one wall is $2.4$ -- an aggressive shrinkage, not a good
fit.  Our prior has no such fallback.  Its scale is fixed by the nominal
coefficients, so where the data cannot pin the field down it keeps extrapolating
with the confidence the equation gave it, and produces something worse than
saying nothing.

The practical consequence is a deployment rule we can state and a reviewer can
check: this construction should be paired with a fallback that detects when the
posterior is extrapolating beyond what the observations constrain and reverts to
shrinkage.  We do not have that mechanism, and until it exists the honest scope
of the method is the layouts where it ties the oracle rather than the ones where
it might have helped most.

We report the prediction and its refutation in this form because three earlier
mechanism predictions in this project were also wrong, and each was easy to
re-narrate afterwards.  A prediction that cannot fail is not evidence, and the
isotropic control is what made this one able to fail -- and what stopped us from
attributing the result to anisotropy after the fact.
